\section{Abstract}
\label{Abstract}
\setlength{\parindent}{20pt}
\justifying

\noindent
While previous bilingual statistical learning experiments found varied bilingual advantages, these were always achieved in the context of an explicit language cue for participants to track regularities separately in both languages. The present work constructs a novel bilingual statistical learning task in which this cue is implicit, and must be implied from monitoring the co-occurrence of morphemes. Over the course of three studies, the paradigm is refined and simplified, and finds evidence for group-level segmentation and co-occurrence tracking to differentiate novel between-language from within-language morphemic pairings. Participants were from a very wide range of language backgrounds, and data for up to four of their spoken languages was used to form multiple continuous predictors. Analysis of individual differences was aided by the development of a more exact language balance measure, which was responsible for the finding of possible quadrilingual difficulties in co-occurrence tracking.